\begin{proof}[Proof of Theorem11]
We use a similar construction as in Theorem10. The difference is that now we want a problem in $\mathsf{BPP}$ that is infinitely-often hard, and a problem in $\mathsf{UP}$ that is almost-everywhere hard, therefore we need to define the oracles differently. In particular, the oracle $P$ only receives one input $r$, and the oracle $Q$ receives two inputs $(x,r)$, where $|r|=10|x|$. The hard language in $\mathsf{BPP}$ is:
\[L_P=\{1^n:\Pr_{r\sim\{0,1\}^{n^4}}[P(r)=1]\geq1/2\}.\]
It is guaranteed that for every $n$, the fraction of length-$n^4$ strings that are in $P$ is either at most $1/3$ or at least $2/3$, therefore $L_P\in\mathsf{BPP}$. On the other hand, the hard language in $\mathsf{UP}$ is:
\[L_Q=\{x:\exists r\in\{0,1\}^{10|x|},Q(x,r)=1\}.\]
It is guaranteed that for every input $x$, there is at most one string $r\in\{0,1\}^{10|x|}$ such that $Q(x,r)=1$, therefore $L_Q\in\mathsf{UP}$. We will construct the oracles $P$ and $Q$ such that $L_P\notin\mathsf{DTIME}[2^n]$ and $L_Q\notin\mathrm{i.o.}\,\mathsf{BPTIME}[2^n]$. Note that $L_P$ is a unary language, therefore we also have $\mathsf{BPP}\not\subseteq\mathrm{heur}\text{-}\mathsf{DTIME}[2^n]$.
We use the same terminologies as before: $M_1,M_2,\ldots$ is an enumeration of deterministic Turing machines running in $2^n$ time, and $N_1,N_2,\ldots$ is an enumeration of probabilistic Turing machines running in $2^n$ time. The $\ell$-th slice of $P$ is the set of inputs with length $\ell$, and the $n$-th slice of $Q$ is the set of inputs of the form $(x,r)$ where $|x|=n$ and $|r|=10n$. An entry is fixed if its value will never change in the subsequent construction. We choose $n_1$ to be a large enough constant and define a sequence $\{n_i\}$ where $n_i=2^{10n_{i-1}}$ for each $i\geq2$. We want every $N_i$ to fail to solve $L_Q$ on all but finitely many input lengths, and every $M_i$ to fail to solve $L_P$ on input $1^{n_i}$.
Our construction proceeds in stages. We guarantee that:
\begin{itemize}
\item For every input length $n$ that is not in the sequence $\{n_i\}$, the $n$-th slice of $P$ is identically $0$.
\item For every input length $n$, there are at most $2n$ entries in the entire $n$-th slice of $Q$ that returns $1$.
\item Consider only the oracle $Q$. Before the $n$-th stage, the number of fixed entries beyond the $n$-th slice of $Q$ (including the $n$-th one) is at most $2^{6n}$, and all of them are fixed to be $0$.
\end{itemize}
We start from the $n^{\star}$-th stage for some large enough constant $n^{\star}$. For each $n\geq n^{\star}$, in the $n$-th stage, we diagonalize against the machines $N_i$ and $M_i$, and fix the $n$-th slice of $Q$ as well as the $n^4$-th slice of $P$ (when $n$ equals some $n_i$). Details follow.
\paragraph{Diagonalization against $N_i$.}
Let $n'$ be the smallest integer in the sequence $\{n_i\}$ such that $n'\geq n$, and $i'$ be the index such that $n'=n_{i'}$. The $(n')^4$-th slice of $P$ may contain either few or most strings, and we do not know which is the case yet. Therefore, we will diagonalize each $N_i$ twice, first assuming the $n'$-th slice of $P$ is nearly empty and then assuming the $n'$-th slice of $P$ is nearly full. For this reason, we will need $2n$ distinct strings $x_{i,b}\in\{0,1\}^n$, one for each pair of $i\in[n]$ and $b\in\{0,1\}$. Note that at the beginning of the $n$-th stage, there are at most $2^{6n}\ll2^{10n}$ strings fixed in the $n$-th slice of $Q$, therefore we can always choose $2n$ strings $x_{i,b}$ where no entry of the form $Q(x_{i,b},r)$ is fixed.
For each $i\in[n]$ and $b\in\{0,1\}$, we simulate $N_i(x_{i,b})$, assuming that most entries in the $n'$-th slice of $P$ returns $b$. Note that $N_i$ is a probabilistic machine running in time $2^n$, thus it has $B:=2^{2^n}$ computational branches, where each branch has $1/B$ probability mass. On each branch, whenever $N_i$ asks a query $P(r)$ or $Q(x,r)$, if this query is already fixed, then we return the corresponding value fixed before; otherwise:
\begin{itemize}
\item Suppose the query is $P(r)$. If $|r|$ is in the sequence $\{n_i^4\}$, then since $P(r)$ is not fixed and $|r|\leq2^n$, we know that $|r|=(n')^4$ and we return $b$ as the answer. Otherwise ($|r|\neq(n')^4$) we return $0$.
\item Suppose the query is $Q(x,r)$, then we return $0$ as the answer.
\end{itemize}
Again, we can only afford to fix the heavy entries, defined as follows. Let $K$ be the number of machine-input pairs that we still need to simulate in the future before fixing the $(n')^4$-th slice of $P$; note that this includes both the current $N_i(x_{i,b})$ and the final $M_{i'}(1^{n'})$. That is:
\[K=2(n-i)+(2-b)+\sum_{j=n+1}^{n'}(2j)+1.\]
\begin{itemize}
\item An entry $P(r)$ is heavy if $|r|=(n')^4$ and at least $2^{-(n'+4)K}$ fraction of branches queried $P(r)$.
\item An entry $Q(x,r)$ is heavy if at least $\frac{1}{10^4|x|^3}$ fraction of branches queried $Q(x,r)$.
\end{itemize}
We fix every heavy entry according to the way we answer this entry before. That is, every heavy entry $P(r)$ is fixed to be $b$ and every heavy entry $Q(x,r)$ is fixed to be $0$. The number of heavy entries in $P$ is at most $2^{(n'+4)K}\cdot2^n$; for each input length $m\geq n$, the number of heavy entries in the $m$-th slice of $Q$ is at most $2^n\cdot10^4\cdot m^3$.
Let $p$ be the probability (over the $B$ computational branches) that $N_i(x_{i,b})$ outputs $1$. If $p<1/2$ then we find a string $r$ such that $Q(x_{i,b},r)$ is unfixed yet and set $Q(x_{i,b},r)=1$, making $x_{i,b}\in L_Q$. Such a string $r$ exists since we have fixed strictly less than $2^{10n}$ strings in the $n$-th slice of $Q$. If $p\geq1/2$ then we do nothing.
Assuming that most entries in the $(n')^4$-th slice of $P$ indeed returns $b$, at this point the probability that $N_i$ outputs the correct answer on input $x_{i,b}$ is at most $1/2$. However, the oracles $P$ and $Q$ might change in the future, thus the behavior of $N_i$ might also change. What fraction of computational branches might be affected in the future?
\begin{itemize}
\item Consider the $K-1$ machine-input pairs that we will simulate in the future before the $(n')^4$-th slice of $P$ is completely fixed. The number of heavy entries in the $(n')^4$-th slice of $P$ that these machines fix is at most
\[\sum_{k=1}^{K-1}2^{(n'+4)k}\cdot2^{n'}\leq2\cdot2^{(n'+4)(K-1)}\cdot2^{n'}.\]
Therefore, the probability mass of computational branches of $N_i(x_{i,b})$ influenced by the future changes to the $(n')^4$-th slice of $P$ is at most
\[2\cdot2^{(n'+4)(K-1)}\cdot2^{n'}\cdot2^{-(n'+4)K}\leq1/8.\]
\item For each $m\leq2^n$, there are at most $2m$ entries in the $m$-th slice of $Q$ that we will fix to be $1$. The probability mass of computational branches of $N_i(x_{i,b})$ influenced by these entries is at most
\[\sum_{m=n}^{2^n}\frac{2m}{10^4m^3}<\frac{1}{10^3}.\]
\end{itemize}
It follows that the future modifications to the oracles $P$ and $Q$ will only affect the accept probability of $N_i(x_{i,b})$ by $\frac18+\frac{1}{10^3}<\frac16$. Therefore, it cannot be the case that
\[\Pr[N_i(x_{i,b})=L_Q(x_{i,b})]\geq2/3.\]
\paragraph{Diagonalization against $M_i$.}
If $n=n_i$ for some $i\in\mathbb N$, then we want that $M_i$ does not compute $L_P$ on input $1^n$ correctly. We simulate $M_i$ on the input $1^n$. When $M_i$ asks a query $P(r)$ or $Q(x,r)$, if this query is already fixed, then we return the corresponding value fixed before; otherwise we return $0$ and fix this queried entry.
After this simulation, we argue that only a small fraction of entries in the $n^4$-th slice of $P$ are fixed. First, the simulation of $M_i(1^n)$ fixes at most $2^n$ entries. Second, there are at most $K=\sum_{j=1}^n(2j)=O(n^2)$ machine-input pairs $N_i(x_{i,b})$ simulated so far, and the $k$-th such machine fixes at most $2^{(n+4)k}\cdot2^n$ entries. Therefore, the number of entries fixed by some previous $N_i(x_{i,b})$ is at most
\[\sum_{k=1}^K2^{(n+4)k}\cdot2^n\leq2^{O(n^3)}.\]
It follows that all but $2^{O(n^3)}$ entries in the $n^4$-th slice of $P$ are not fixed yet. Let $b\in\{0,1\}$ be the output bit of $M_i(1^n)$, then we set $L_P(1^n)=1-b$ by fixing every unfixed entry in the $n^4$-th slice of $P$ to be $1-b$. We have that on input $1^n$, $L_P$ satisfies the $\mathsf{BPP}$ promise, and $M_i$ fails to solve $L_P$.
\paragraph{Clear-up.}
At the end of stage $n$, we fix every unfixed input $Q(x,r)$ on the $n$-th slice to be $0$. It is easy to see that there are at most $2n$ entries in the $n$-th slice of $Q(x,r)$ that returns $1$.
During the $n$-th stage, the number of entries we fixed beyond the $(n+1)$-st slice of $Q$ is at most
\[2^n+\sum_{m=n+1}^{2^n}2^n\cdot10^4\cdot m^3\leq2^{5n}.\]
Thus the total number of entries we fixed beyond the $(n+1)$-st slice of $Q$ is at most $2^{6n}+2^{5n}<2^{6(n+1)}$.
\end{proof}
